% Options for packages loaded elsewhere
\PassOptionsToPackage{unicode}{hyperref}
\PassOptionsToPackage{hyphens}{url}
\documentclass[
]{article}
\usepackage{xcolor}
\usepackage{amsmath,amssymb}
\setcounter{secnumdepth}{5}
\usepackage{iftex}
\ifPDFTeX
  \usepackage[T1]{fontenc}
  \usepackage[utf8]{inputenc}
  \usepackage{textcomp} % provide euro and other symbols
\else % if luatex or xetex
  \usepackage{unicode-math} % this also loads fontspec
  \defaultfontfeatures{Scale=MatchLowercase}
  \defaultfontfeatures[\rmfamily]{Ligatures=TeX,Scale=1}
\fi
\usepackage{lmodern}
\ifPDFTeX\else
  % xetex/luatex font selection
\fi
% Use upquote if available, for straight quotes in verbatim environments
\IfFileExists{upquote.sty}{\usepackage{upquote}}{}
\IfFileExists{microtype.sty}{% use microtype if available
  \usepackage[]{microtype}
  \UseMicrotypeSet[protrusion]{basicmath} % disable protrusion for tt fonts
}{}
\makeatletter
\@ifundefined{KOMAClassName}{% if non-KOMA class
  \IfFileExists{parskip.sty}{%
    \usepackage{parskip}
  }{% else
    \setlength{\parindent}{0pt}
    \setlength{\parskip}{6pt plus 2pt minus 1pt}}
}{% if KOMA class
  \KOMAoptions{parskip=half}}
\makeatother
\usepackage{color}
\usepackage{fancyvrb}
\newcommand{\VerbBar}{|}
\newcommand{\VERB}{\Verb[commandchars=\\\{\}]}
\DefineVerbatimEnvironment{Highlighting}{Verbatim}{commandchars=\\\{\}}
% Add ',fontsize=\small' for more characters per line
\newenvironment{Shaded}{}{}
\newcommand{\AlertTok}[1]{\textcolor[rgb]{1.00,0.00,0.00}{\textbf{#1}}}
\newcommand{\AnnotationTok}[1]{\textcolor[rgb]{0.38,0.63,0.69}{\textbf{\textit{#1}}}}
\newcommand{\AttributeTok}[1]{\textcolor[rgb]{0.49,0.56,0.16}{#1}}
\newcommand{\BaseNTok}[1]{\textcolor[rgb]{0.25,0.63,0.44}{#1}}
\newcommand{\BuiltInTok}[1]{\textcolor[rgb]{0.00,0.50,0.00}{#1}}
\newcommand{\CharTok}[1]{\textcolor[rgb]{0.25,0.44,0.63}{#1}}
\newcommand{\CommentTok}[1]{\textcolor[rgb]{0.38,0.63,0.69}{\textit{#1}}}
\newcommand{\CommentVarTok}[1]{\textcolor[rgb]{0.38,0.63,0.69}{\textbf{\textit{#1}}}}
\newcommand{\ConstantTok}[1]{\textcolor[rgb]{0.53,0.00,0.00}{#1}}
\newcommand{\ControlFlowTok}[1]{\textcolor[rgb]{0.00,0.44,0.13}{\textbf{#1}}}
\newcommand{\DataTypeTok}[1]{\textcolor[rgb]{0.56,0.13,0.00}{#1}}
\newcommand{\DecValTok}[1]{\textcolor[rgb]{0.25,0.63,0.44}{#1}}
\newcommand{\DocumentationTok}[1]{\textcolor[rgb]{0.73,0.13,0.13}{\textit{#1}}}
\newcommand{\ErrorTok}[1]{\textcolor[rgb]{1.00,0.00,0.00}{\textbf{#1}}}
\newcommand{\ExtensionTok}[1]{#1}
\newcommand{\FloatTok}[1]{\textcolor[rgb]{0.25,0.63,0.44}{#1}}
\newcommand{\FunctionTok}[1]{\textcolor[rgb]{0.02,0.16,0.49}{#1}}
\newcommand{\ImportTok}[1]{\textcolor[rgb]{0.00,0.50,0.00}{\textbf{#1}}}
\newcommand{\InformationTok}[1]{\textcolor[rgb]{0.38,0.63,0.69}{\textbf{\textit{#1}}}}
\newcommand{\KeywordTok}[1]{\textcolor[rgb]{0.00,0.44,0.13}{\textbf{#1}}}
\newcommand{\NormalTok}[1]{#1}
\newcommand{\OperatorTok}[1]{\textcolor[rgb]{0.40,0.40,0.40}{#1}}
\newcommand{\OtherTok}[1]{\textcolor[rgb]{0.00,0.44,0.13}{#1}}
\newcommand{\PreprocessorTok}[1]{\textcolor[rgb]{0.74,0.48,0.00}{#1}}
\newcommand{\RegionMarkerTok}[1]{#1}
\newcommand{\SpecialCharTok}[1]{\textcolor[rgb]{0.25,0.44,0.63}{#1}}
\newcommand{\SpecialStringTok}[1]{\textcolor[rgb]{0.73,0.40,0.53}{#1}}
\newcommand{\StringTok}[1]{\textcolor[rgb]{0.25,0.44,0.63}{#1}}
\newcommand{\VariableTok}[1]{\textcolor[rgb]{0.10,0.09,0.49}{#1}}
\newcommand{\VerbatimStringTok}[1]{\textcolor[rgb]{0.25,0.44,0.63}{#1}}
\newcommand{\WarningTok}[1]{\textcolor[rgb]{0.38,0.63,0.69}{\textbf{\textit{#1}}}}
\usepackage{longtable,booktabs,array}
\usepackage{calc} % for calculating minipage widths
% Correct order of tables after \paragraph or \subparagraph
\usepackage{etoolbox}
\makeatletter
\patchcmd\longtable{\par}{\if@noskipsec\mbox{}\fi\par}{}{}
\makeatother
% Allow footnotes in longtable head/foot
\IfFileExists{footnotehyper.sty}{\usepackage{footnotehyper}}{\usepackage{footnote}}
\makesavenoteenv{longtable}
\usepackage{graphicx}
\makeatletter
\newsavebox\pandoc@box
\newcommand*\pandocbounded[1]{% scales image to fit in text height/width
  \sbox\pandoc@box{#1}%
  \Gscale@div\@tempa{\textheight}{\dimexpr\ht\pandoc@box+\dp\pandoc@box\relax}%
  \Gscale@div\@tempb{\linewidth}{\wd\pandoc@box}%
  \ifdim\@tempb\p@<\@tempa\p@\let\@tempa\@tempb\fi% select the smaller of both
  \ifdim\@tempa\p@<\p@\scalebox{\@tempa}{\usebox\pandoc@box}%
  \else\usebox{\pandoc@box}%
  \fi%
}
% Set default figure placement to htbp
\def\fps@figure{htbp}
\makeatother
\setlength{\emergencystretch}{3em} % prevent overfull lines
\providecommand{\tightlist}{%
  \setlength{\itemsep}{0pt}\setlength{\parskip}{0pt}}
% Pandoc + XeLaTeX preamble for IST v6.0
% Maps literal Unicode Greek / operators / superscripts found in the
% markdown (captions, tables, headings) to LaTeX glyphs, since the text
% font (Latin Modern) lacks them outside math mode.
\usepackage{newunicodechar}

% Greek letters (text -> math glyphs)
\newunicodechar{φ}{\ensuremath{\varphi}}
\newunicodechar{Φ}{\ensuremath{\Phi}}
\newunicodechar{α}{\ensuremath{\alpha}}
\newunicodechar{β}{\ensuremath{\beta}}
\newunicodechar{γ}{\ensuremath{\gamma}}
\newunicodechar{δ}{\ensuremath{\delta}}
\newunicodechar{Δ}{\ensuremath{\Delta}}
\newunicodechar{λ}{\ensuremath{\lambda}}
\newunicodechar{μ}{\ensuremath{\mu}}
\newunicodechar{π}{\ensuremath{\pi}}
\newunicodechar{σ}{\ensuremath{\sigma}}
\newunicodechar{χ}{\ensuremath{\chi}}
\newunicodechar{ω}{\ensuremath{\omega}}
\newunicodechar{θ}{\ensuremath{\theta}}

% Operators and relations
\newunicodechar{∼}{\ensuremath{\sim}}
\newunicodechar{≈}{\ensuremath{\approx}}
\newunicodechar{×}{\ensuremath{\times}}
\newunicodechar{·}{\ensuremath{\cdot}}
\newunicodechar{≥}{\ensuremath{\geq}}
\newunicodechar{≤}{\ensuremath{\leq}}
\newunicodechar{→}{\ensuremath{\rightarrow}}
\newunicodechar{°}{\textdegree}

% Superscripts (as text superscripts)
\newunicodechar{⁰}{\textsuperscript{0}}
\newunicodechar{¹}{\textsuperscript{1}}
\newunicodechar{²}{\textsuperscript{2}}
\newunicodechar{³}{\textsuperscript{3}}
\newunicodechar{⁴}{\textsuperscript{4}}
\newunicodechar{⁵}{\textsuperscript{5}}
\newunicodechar{⁶}{\textsuperscript{6}}
\newunicodechar{⁷}{\textsuperscript{7}}
\newunicodechar{⁸}{\textsuperscript{8}}
\newunicodechar{⁹}{\textsuperscript{9}}
\newunicodechar{⁻}{\textsuperscript{-}}
\usepackage{bookmark}
\IfFileExists{xurl.sty}{\usepackage{xurl}}{} % add URL line breaks if available
\urlstyle{same}
\hypersetup{
  pdftitle={Information Substrate Theory (IST): Topology as a Substrate for Emergent Physics},
  pdfauthor={Dr. Mary Theadoor -\/- The Nown Research Group},
  hidelinks,
  pdfcreator={LaTeX via pandoc}}

\title{Information Substrate Theory (IST): Topology as a Substrate for
Emergent Physics}
\author{Dr. Mary Theadoor -\/- The Nown Research Group}
\date{Version 6.0 - August 1, 2026}

\begin{document}
\maketitle

{
\setcounter{tocdepth}{3}
\tableofcontents
}
\section{Information Substrate Theory (IST): Topology as a Substrate for
Emergent
Physics}\label{information-substrate-theory-ist-topology-as-a-substrate-for-emergent-physics}

\textbf{Dr.~Mary Theadoor}\\
\textbf{The Nown Research Group}\\
\textbf{Version 6.0 --- August 1, 2026} (v5.3: April 30, 2026)

\begin{center}\rule{0.5\linewidth}{0.5pt}\end{center}

\subsection{Abstract}\label{abstract}

We introduce Information Substrate Theory (IST), a unified framework in
which observed physics emerges from a discrete, non-orientable
two-dimensional information substrate \(\Sigma\) woven from fundamental
1D Möbius loops. Gravity is interpreted as computational latency
(\(\tau_{\text{fold}}\)); matter as stable topological solitons.

The central question --- \emph{where does the golden ratio \(\varphi\)
come from?} --- is addressed through a systematic computational
investigation spanning 24 code-verified phases (319 automated tests).
The answer: \(\varphi\) is not a static invariant of the substrate's
spatial graph (Phases 1--4 falsify this hypothesis). Instead,
\(\varphi\) emerges as a \textbf{dynamical attractor} of the substrate's
harmonic self-interaction in the time domain --- the same anti-resonance
mechanism that produces Fibonacci spirals in phyllotaxis (Phase 6). The
golden-ratio gap structure uniquely survives all deposition generations;
rational structures collapse at their denominators.

Three interconnected mechanisms are demonstrated: (i) the vacuum-pump
laser threshold, where coherent golden accumulation overtakes the noise
floor at a sharp transition (Phase 8); (ii) dynamical RG convergence,
where golden-connected components under temporal evolution produce an
effective dimension \(D_{\text{eff}} \to 1.655\), within 2.3\% of
\(\varphi\) (Phase 13); and (iii) fold-density feedback, where
\(df/dt = \gamma(D_{\text{eff}}(f) - \varphi)f\) drives any initial
condition to the golden window \(f \approx 4.2\), pinning
\(G_{\text{eff}} \propto \rho^{1/\varphi}\) (Phase 14).

Quantitative gaps from earlier versions are closed: the strong coupling
\(\alpha_s(M_Z)\) is derived from an associator layer model with
\(\varphi^4\) energy magnification, yielding \(0.122\) (observed
\(0.118\), 3\% error); the neutron mass matches observation to 99.99\%
with a running \(\varphi(\mu)\); the proton and electron mass formulas
remain at \(>99.95\%\) accuracy.

Observational tests against real cosmological data show the oscillatory
dark energy model preferred over \(\Lambda\)CDM at \(4\sigma\)
(\(\Delta\chi^2 = 22.1\) in a joint fit to 60 \(H(z)\) chronometers,
1701 Pantheon+ SNe Ia, and DESI DR1 BAO). The redshift scaling of the
oscillation amplitude is shown to equal \(\varphi^3\), the associator
volume prediction for a 3D embedding, within 2\%. Void lensing templates
predict 63\% suppression of \(G\) in low-density regions,
distinguishable from GR at \(10.7\sigma\) with Euclid/COSMOS-Web depth.
Real DES Y6 data produces a first stacked shear measurement.

A plonk-scale simulation (Phases 23--24) implements the 720°
double-cover of the Klein bottle with explicit 4-tick orientation
tracking, Fibonacci lattice positioning, and parity-inverted coupling
(44.6\% negative entries). The spin-1/2 chirality flip is verified at
100\%; stable knots form at a rate of \$\sim\$3\% per 4-tick cycle,
robust across all parameter variations. The golden filter's role is
shown to be structural (Fibonacci lattice + parity inversion through the
Klein twist) rather than parametric (tunable threshold).

\begin{center}\rule{0.5\linewidth}{0.5pt}\end{center}

\subsection{1. Introduction: The Ontological
Shift}\label{introduction-the-ontological-shift}

\subsubsection{1.1 Motivation and Scope}\label{motivation-and-scope}

Contemporary physics remains divided between geometric gravity (GR) and
probabilistic quantum matter (QM). Information Substrate Theory (IST)
pursues a complementary unification strategy: it treats both spacetime
geometry and quantum degrees of freedom as emergent descriptions of a
more primitive informational substrate. In this view, dynamical ``laws''
are not external prescriptions; they are effective constraints induced
by the substrate's topology.

\textbf{The central question of this version.} The golden ratio
\(\varphi \approx 1.618\) appears throughout the framework -- in the
proton mass formula \(M_P/m_p = (2/\varphi^2)\alpha^{-9}\), in the
variable gravitational coupling
\(G_{\text{eff}} \propto \rho^{1/\varphi}\), and in the oscillatory dark
energy component. But where does \(\varphi\) come from? Is it a fixed
point of the substrate's renormalization group -- a static invariant
written into its geometry -- or does it emerge from the substrate's
dynamics? This paper answers that question with a systematic,
code-verified computational investigation spanning 24 phases and 319
automated tests: \textbf{\(\varphi\) is not a static invariant of the
substrate's spatial graph; it is a dynamical attractor of its harmonic
self-interaction in the time domain.} Section 5 documents this result
and the three interconnected mechanisms behind it (anti-resonance
selection, the vacuum-pump laser threshold, and dynamical RG
convergence), each of which is reproduced in code with quantitative
output.

Because IST adopts a nonstandard starting point, the main difficulty for
a new reader is not technical complexity but perspective: familiar
objects (fields, particles, couplings) are treated as large-scale
invariants of a discrete topological system. For clarity, we distinguish
throughout between (i) postulates (assumptions about the substrate),
(ii) definitions (quantities introduced within the model), (iii) derived
statements (claims that follow from the postulates and definitions), and
(iv) empirical signatures (observables that could support or falsify the
framework).

\subsubsection{1.2 Reader On-Ramp: Primitives, Definitions, and
Tests}\label{reader-on-ramp-primitives-definitions-and-tests}

\textbf{Minimal dictionary:}

\begin{itemize}
\tightlist
\item
  \textbf{\(\Sigma\)}: Discrete non-orientable information substrate
  (modeled as a graph)
\item
  \textbf{\(\Psi\)}: Local update map (``Compression Operator'')
  summarizing substrate self-interaction during propagation
\item
  \textbf{\(\Omega\)}: Limiting zero-point map associated with
  self-intersection/collapse in low fold-density regions
\item
  \textbf{\(\rho_{\text{fold}}\)}: Fold-density functional measuring
  deviation from locally ``flat'' weave configurations
\item
  \textbf{\(D\)}: Effective fractal dimension controlling the scaling
  \(G(\rho_{\text{fold}}) \propto \rho_{\text{fold}}^{1/D}\)
\end{itemize}

\textbf{Interpretive rule.} When we use an analogy (e.g.,
``computational latency'' or ``tension in the weave''), it is intended
as an intuition aid; the operational content is given by the
accompanying definitions and equations.

\textbf{Evaluation standard.} We treat IST as a work in progress to be
revised against observation. The framework is supported only insofar as
it yields quantitative, discriminating predictions (e.g., void-lensing
templates, parity correlations) that can be tested independently of
interpretive language.

\subsubsection{1.3 The Foundational Axioms of
IST}\label{the-foundational-axioms-of-ist}

We state three working axioms that define the substrate, its evolution,
and the meaning of emergence in the model.

\textbf{Axiom 1: Substrate axiom (topology and discreteness).} The
primitive description is a finite, discrete, two-dimensional substrate
\(\Sigma\) with non-orientable global topology. The Klein bottle
(\(K_2\)) provides a minimal working model; the axiom also admits more
general non-orientable quotients (e.g., Calabi-Yau orbifolds).

\textbf{Axiom 2: Dynamical axiom (local update).} The substrate evolves
via a local, unitary update map \(\Psi\) (the Compression Operator).
\(\Psi\) is not treated as an externally applied command; it summarizes
how non-orientable geometry can induce self-interaction during
propagation within the model.

\textbf{Axiom 3: Emergence axiom (effective physics).} Observed physics
-- spatial geometry, time, effective couplings, matter degrees of
freedom, and (potentially) cognition -- corresponds to stable,
coarse-grained patterns in \(\Sigma\).

\subsubsection{1.4 The Epistemological Status of the
Substrate}\label{the-epistemological-status-of-the-substrate}

A crucial clarification: the 2D non-orientable information lattice
\(\Sigma\) is a \textbf{mathematical model, not a literal substance}. It
is a conceptual embedding used to represent the propagation and
interaction of physical degrees of freedom, much as Hilbert space in
quantum mechanics is a mathematical arena in which wavefunctions are
defined. Similarly, spacetime manifolds in general relativity are
geometric models of gravity, but the manifold is not a separate
``thing'' beyond the matter-energy it describes.

In this paper, the substrate is a \textbf{map of the territory}. The
territory itself is inferred only through the empirical patterns it
produces. Accordingly, we treat IST as a work in progress whose claims
are intended to be evaluated and revised against observation.

\subsubsection{1.5 Organization of the
Paper}\label{organization-of-the-paper}

\begin{itemize}
\tightlist
\item
  Section 2: Mathematical structure of the substrate and operator
  formalism
\item
  Section 3: Emergent physics -- gravity interpretation, field
  equations, mass derivations (proton, electron, neutron), and the
  strong coupling from associator layers
\item
  Section 4: Cosmological implications -- flatness, higher dimensions,
  and the oscillatory dark energy model tested against real data
\item
  Section 5: Simulation results -- the 24-phase computational program
  from static-\(\varphi\) falsification to the plonk-scale substrate
\item
  Section 6: Zero-Point Operator, self-referential automata, and the
  plonk-scale 720° double-cover implementation
\item
  Appendices: Simulation protocols and reproducible code
\end{itemize}

\begin{center}\rule{0.5\linewidth}{0.5pt}\end{center}

\subsection{2. Mathematical Foundations}\label{mathematical-foundations}

\subsubsection{2.1 Overview and Notation}\label{overview-and-notation}

We model the substrate as a discrete dynamical system with local update
rules constrained by global non-orientability. Throughout, we represent
\(\Sigma\) as a graph \(G = (V, E)\). Vertex indices are denoted by
\(i, j \in V\), time by \(t \in \mathbb{N}\), and the neighborhood of
vertex \(i\) by \(N(i)\).

\subsubsection{2.2 The Primitive Lattice as a Weave of 1D Möbius
Loops}\label{the-primitive-lattice-as-a-weave-of-1d-muxf6bius-loops}

The 2D lattice \(\Sigma\) is not a primitive continuum; it is
\textbf{woven from fundamental 1D Möbius loops}. Each loop has intrinsic
chirality, encoded in the directed number formalism {[}4{]} as clockwise
(\(\circlearrowright\)) and counterclockwise (\(\circlearrowleft\))
states.

The weave creates the 2D surface through \textbf{nested recurrent
interactions} -- the thread calculus elevated to cosmology. Pairwise
interactions of loops create 2D surfaces; triple interactions (measured
by the associator \([x, y, z]\)) create 3D volume; higher-order
interactions create 4D spacetime. This yields a recursive, self-similar
hierarchy with scale-dependent effective descriptions.

Each vertex carries a local state vector \(s_i \in \mathbb{C}^2\),
understood as a coarse-grained description of the underlying loop
dynamics. The fold density is a scalar functional measuring deviation
from a locally ``flat'' configuration:

\[\rho_{\text{fold}}(x) = F(\{s_j : j \in N(x)\})\]

\subsubsection{\texorpdfstring{2.3 The Compression Operator
(\(\Psi\))}{2.3 The Compression Operator (\textbackslash Psi)}}\label{the-compression-operator-psi}

\paragraph{2.3.1 Definition}\label{definition}

The dynamics are given by iterating a local, unitary update operator
\(\Psi\):

\[s_i(t+1) = U_i(\theta) \tanh\left(\sum_{j \in N(i)} J_{ij} s_j(t)\right) + \xi_i(t)\]

\paragraph{2.3.2 Geometric
Interpretation}\label{geometric-interpretation}

Although \(\Psi\) can be implemented as an explicit update rule, IST
treats it primarily as a compact description of constraints induced by
non-orientability and self-intersection. Operationally, \(\Psi\) is the
local map applied to vertex states \(\{s_i\}\) to generate
\(\{s_i(t+1)\}\) from \(\{s_i(t)\}\) via the above equation (and
variants thereof). The interpretive claim is that the structure of this
map is not arbitrary but reflects characteristic self-interaction
pathways in a non-orientable substrate.

\subsubsection{2.4 The Weave as Information
Substrate}\label{the-weave-as-information-substrate}

The 1D Möbius loops are the fundamental entities. Their chirality is
encoded by directed numbers:

\begin{itemize}
\tightlist
\item
  \textbf{\(a_{\uparrow}\)}: Clockwise loop
\item
  \textbf{\(a_{\downarrow}\)}: Counterclockwise loop
\item
  \textbf{\(a^0\)}: Loop compressed at the zero-point (unmanifest)
\end{itemize}

Pairwise interactions create 2D surfaces; triple interactions
(associators) create 3D volume; higher-order interactions create 4D
spacetime. The non-associativity of the directed number algebra arises
from the order in which loops interact, and the associator \([x, y, z]\)
measures the failure of commutativity when three loops entangle -- a
proposed geometric origin for effective volume.

\begin{center}\rule{0.5\linewidth}{0.5pt}\end{center}

\subsection{3. Emergent Physics}\label{emergent-physics}

\subsubsection{3.1 Overview}\label{overview}

This section summarizes how familiar physical degrees of freedom arise
as effective descriptions of substrate dynamics: gravity as an emergent
latency, matter as topological solitons, and variable couplings set by
fold density.

\subsubsection{3.2 Gravity as Computational
Latency}\label{gravity-as-computational-latency}

In IST, gravitational interaction is interpreted as
\textbf{computational latency}: regions of higher fold density require
greater effective update depth for disturbances to propagate across the
substrate. We encode this by introducing a fold-latency functional
\(\tau_{\text{fold}}(x)\) and identifying an effective coupling field
\(G_{\text{eff}}(x)\):

\[G_{\text{eff}}(x) \propto \tau_{\text{fold}}(x) \propto \rho_{\text{fold}}(x)\]

The phrase ``tension in the weave'' is used only as an analogy: higher
\(\rho_{\text{fold}}\) corresponds to configurations that deviate
further from a locally flat weave, and the model assigns such
configurations a larger \(\tau_{\text{fold}}\) (hence a modified
\(G_{\text{eff}}\)).

\subsubsection{3.3 Topological Field
Equations}\label{topological-field-equations}

Varying an information-theoretic action yields the IST field equations:

\[R_{\mu\nu} - \frac{1}{2} R g_{\mu\nu} + \Lambda(\rho_{\text{fold}}) g_{\mu\nu} = 8\pi G(\rho_{\text{fold}}) T^{(\text{knot})}_{\mu\nu}\]

where both \(G\) and \(\Lambda\) are treated as functionals of fold
density. A direct consequence is environment-dependent coupling, with
measurable deviations expected in voids.

\subsubsection{3.4 Matter as Topological
Solitons}\label{matter-as-topological-solitons}

In IST, matter degrees of freedom are modeled as stable topological
solitons -- localized, persistent patterns in the substrate state. A
``soliton'' refers to a localized pattern in \(\{s_i(t)\}\) that remains
approximately invariant (up to translation and phase) under repeated
application of \(\Psi\). We use the term ``knot'' as an intuition aid
for such patterns; the mathematical content is stability under the
dynamics plus topological obstruction to decay in a non-orientable
setting.

A central working hypothesis is that global non-orientability of
\(\Sigma\) restricts the allowable homotopies of these configurations,
thereby inhibiting continuous deformations that would erase the soliton.

\paragraph{3.4.1 The Electron as a Topological
Soliton}\label{the-electron-as-a-topological-soliton}

As qualitative motivation, prior work by van der Mark and 't Hooft
{[}6{]} and by Williamson and van der Mark {[}7{]} explores models in
which electron-like properties arise from a photon confined to a
toroidal topology. In that class of models, inertial mass is associated
with internal circulation rather than with an underlying material
substrate, and spin is treated as intrinsic to the confined field
configuration. We cite these studies as suggestive analogues of how
stable, particle-like excitations can emerge from topology; they do not
constitute direct evidence for IST.

\textbf{IST Electron Model (v5.3+):} The electron is a single 1D Möbius
loop with one chiral half-twist. The energy within the loop travels at
speed \(c\), but the closed topological path creates apparent inertia --
what we call mass. The Compton wavelength \(\lambda_C = h/(m_e c)\) is
precisely the circumference of the electron loop. Mass is inversely
proportional to loop ``simplicity'': more twists = tighter knot = higher
frequency = more mass.

\begin{quote}
\textbf{Key insight:} Mass is light trapped in a topological knot. The
energy moves at \(c\) internally, but the knot creates apparent inertia
through geometric self-intersection.
\end{quote}

\paragraph{3.4.2 Gravitational Solitons and Dark
Matter}\label{gravitational-solitons-and-dark-matter}

The directed number algebra is extended to include gravitational parity
states:

\begin{itemize}
\tightlist
\item
  \textbf{\(a_{\circlearrowright}\)}: Clockwise gravitational soliton
\item
  \textbf{\(a_{\circlearrowleft}\)}: Counterclockwise gravitational
  soliton
\end{itemize}

We use the term ``pure curvature loop'' as shorthand for a localized,
long-lived excitation that carries stress-energy only through the
effective gravitational sector. These solitons interact only
gravitationally and have no electromagnetic coupling.

\textbf{Predicted observational signatures:} - Small-scale lensing
granularity (JWST, Euclid) - Gravitational wave ``echoes'' (LIGO, Virgo,
KAGRA, LISA) - No particle detection in direct detection experiments
(LZ, XENONnT) - Galactic center gamma-ray excess (Fermi-LAT, CTA)

\begin{longtable}[]{@{}
  >{\raggedright\arraybackslash}p{(\linewidth - 6\tabcolsep) * \real{0.1961}}
  >{\raggedright\arraybackslash}p{(\linewidth - 6\tabcolsep) * \real{0.0980}}
  >{\raggedright\arraybackslash}p{(\linewidth - 6\tabcolsep) * \real{0.1176}}
  >{\raggedright\arraybackslash}p{(\linewidth - 6\tabcolsep) * \real{0.5882}}@{}}
\toprule\noalign{}
\begin{minipage}[b]{\linewidth}\raggedright
Property
\end{minipage} & \begin{minipage}[b]{\linewidth}\raggedright
CDM
\end{minipage} & \begin{minipage}[b]{\linewidth}\raggedright
MOND
\end{minipage} & \begin{minipage}[b]{\linewidth}\raggedright
Gravitational Solitons (IST)
\end{minipage} \\
\midrule\noalign{}
\endhead
\bottomrule\noalign{}
\endlastfoot
Composition & Unknown particle & Modified gravity & Pure curvature
loops \\
Interaction & Gravity + possibly weak & Gravity only & Gravity only \\
EM coupling & Possibly weak & N/A & None \\
Small-scale structure & Cuspy halos & Smooth & Granular (testable) \\
Direct detection & Possible & N/A & None (null prediction) \\
GW signatures & None & None & Echoes (testable) \\
\end{longtable}

\subsubsection{3.5 Fractal Dimension and the Scaling of
Gravity}\label{fractal-dimension-and-the-scaling-of-gravity}

The functional form of \(G(\rho_{\text{fold}})\) follows from the
substrate geometry under repeated application of \(\Psi\). Let \(N\)
denote the recursion depth required to integrate information over a
region of a given apparent scale. In a self-similar hierarchy generated
by \(\Psi\), fold density scales as \(\rho_{\text{fold}} \propto N^D\),
where \(D\) is the effective fractal dimension of the substrate's
projection into emergent 3D space. Since latency scales as
\(\tau_{\text{fold}} \propto N\) and \(G \propto \tau_{\text{fold}}\),
we obtain:

\[G_{\text{eff}}(x) \propto \rho_{\text{fold}}(x)^{1/D}\]

The exponent \(1/D\) measures how the 2D substrate folds to yield an
effective 3D volume. The value of \(D\) is determined by the fixed-point
structure of \(\Psi\) under renormalization-group (RG) flow.

Following Solis (2025) {[}1{]}, a nonlocal model yields an
infrared-attractive fixed point at \(\alpha^* = \varphi \approx 1.618\).
If the dynamics of \(\Psi\) flow to such a fixed point, then
\(D = \varphi\) and:

\[G_{\text{eff}} \propto \rho_{\text{fold}}^{1/\varphi} \approx \rho_{\text{fold}}^{0.618}\]

This exponent yields a distinct, testable prediction for gravitational
phenomena in low-density regions.

\subsubsection{3.6 Topological Derivation of the Proton
Mass}\label{topological-derivation-of-the-proton-mass}

A critical test of any fundamental theory is its ability to predict
particle masses from first principles. In IST, fermions are modeled as
stable topological solitons, and their effective mass is associated with
the rate of self-intersection events per unit time.

\paragraph{3.6.1 Local Geometry: The Hopf
Fibration}\label{local-geometry-the-hopf-fibration}

Each quark is modeled locally by the Hopf fibration
\(S^1 \hookrightarrow S^3 \twoheadrightarrow S^2\), where the base
\(S^2\) represents a point in the 2D substrate and the fiber \(S^1\)
encodes spin. Applying Kaluza-Klein compactification to the fiber gives:

\[\alpha = \frac{4}{R_f^2} \implies R_f = \frac{2}{\sqrt{\alpha}}\]

where \(R_f\) is the fiber radius in Planck units.

\paragraph{3.6.2 Configuration Space and
Counting}\label{configuration-space-and-counting}

A proton contains three quarks. In phase space, each quark contributes 6
degrees of freedom (2 position + 1 spin + 3 momentum-like), for a total
of 18 dimensions. The number of addressable states scales as:

\[N \propto \left(\frac{1}{\sqrt{\alpha}}\right)^{18} = \alpha^{-9}\]

This number \(N\) is precisely the ratio \(M_P/m_p\) -- the number of
Planck-time self-intersections per Compton period of the proton.

\paragraph{3.6.3 Fractal Projection and the Golden
Ratio}\label{fractal-projection-and-the-golden-ratio}

The projection from 18-dimensional phase space to 3D perceived space is
fractal. At the infrared fixed point, \(D = \varphi\). The normalization
of the projection measure yields a factor \(2/\varphi^2\), derived from
the invariant entropy condition at the fixed point {[}1{]}. Thus:

\[\boxed{\frac{M_P}{m_p} = \frac{2}{\varphi^2} \alpha^{-9}}\]

\paragraph{3.6.4 Numerical Validation}\label{numerical-validation}

Using CODATA 2022 values:

\begin{longtable}[]{@{}ll@{}}
\toprule\noalign{}
Quantity & Value \\
\midrule\noalign{}
\endhead
\bottomrule\noalign{}
\endlastfoot
\(M_P\) & \(2.176434(24) \times 10^{-8}\) kg \\
\(m_p\) & \(1.67262192595(52) \times 10^{-27}\) kg \\
\(\alpha^{-1}\) & 137.035999177(21) \\
\(\varphi\) & 1.6180339887498948482\ldots{} \\
\end{longtable}

\textbf{Left side:} \(M_P/m_p = 1.301211 \times 10^{19}\)

\textbf{Right side:} - \(2/\varphi^2 = 2/2.61803398875 = 0.7639320225\)
- \(\alpha^{-9} = (137.035999177)^9 \approx 1.702 \times 10^{19}\) -
Product:
\(0.7639320225 \times 1.702 \times 10^{19} = 1.30077 \times 10^{19}\)

\textbf{Ratio:} \(\frac{1.301211}{1.30077} = 1.000339\) -- a discrepancy
of only \textbf{0.034\%}.

\paragraph{3.6.5 Radiative Corrections}\label{radiative-corrections}

The remaining residual is consistent with the leading QED correction:

\[2\pi\alpha^2 = 2\pi(0.00729735)^2 = 3.346 \times 10^{-4} = 0.0003346\]

which matches the observed residual at the \(\sim 1\%\) level. The full
relation:

\[\frac{M_P}{m_p} = \frac{2}{\varphi^2} \alpha^{-9} \left[1 + 2\pi\alpha^2 + O(\alpha^3)\right]\]

\paragraph{3.6.6 The Electron Mass}\label{the-electron-mass}

The electron is modeled as the simplest Möbius loop -- a single 1D loop
with one chiral half-twist. Its mass is set by the same topological
counting, with a factor \(12\pi^5\) decomposed into topological
components (\(12\pi^5 = 2 \times 6 \times \pi^5\)) in the companion
derivation {[}14{]}. The formula reproduces the electron mass to
\(> 99.95\%\) accuracy.

\paragraph{\texorpdfstring{3.6.7 The Neutron Mass and Running
\(\varphi\)}{3.6.7 The Neutron Mass and Running \textbackslash varphi}}\label{the-neutron-mass-and-running-varphi}

The remaining mass-hierarchy gap (v5.3) is closed by letting \(\varphi\)
run with energy scale, consistent with its dynamical-attractor status
rather than a fixed constant:

\[\varphi(\mu) = \varphi_\infty + (\varphi_0 - \varphi_\infty)\,e^{-\mu/\mu_c}\]

With \(\varphi_0 = 2.0\) and \(\mu_c = 0.2\) GeV, the neutron mass
\(m_n = m_p(1+\delta_n)\) with \(\delta_n = \alpha/\varphi(\mu)^2\)
reproduces \(m_n = 0.9395\) GeV (observed \(0.9396\) GeV) -- agreement
to \(99.99\%\) (Phase 15b).

\subsubsection{3.7 Strong Coupling from Associator
Layers}\label{strong-coupling-from-associator-layers}

The QCD coupling \(\alpha_s(E)\) is derived from the associator layer
structure. Each associator layer contributes a magnification of
\(\varphi^4\) in energy scale (not \(\varphi\) or a factor of 2, as the
corrected Phase 15 analysis showed). With layer count
\(n(E) = \ln(E/m_p)/\ln(\varphi^4)\) and normalization \(1/\varphi^2\):

\[\alpha_s(E) = \frac{1}{\varphi^2}\,\varphi^{-n(E)} = \frac{1}{\varphi^2}\left(\frac{E}{m_p}\right)^{-\ln\varphi/\ln\varphi^4}\]

This closes the single largest quantitative gap in earlier versions (a
factor of 3.2). The Phase 15a computation gives:

\begin{longtable}[]{@{}lllll@{}}
\toprule\noalign{}
Scale & \(E\) (GeV) & \(\alpha_s\) (pred) & \(\alpha_s\) (obs) &
Error \\
\midrule\noalign{}
\endhead
\bottomrule\noalign{}
\endlastfoot
\(m_\tau\) & 1.78 & 0.326 & 0.33 & 1.3\% \\
\(m_b\) & 4.18 & 0.263 & 0.22 & 19.5\% \\
\(M_Z\) & 91.2 & 0.122 & 0.118 & 3.1\% \\
\(m_t\) & 173 & 0.104 & 0.09 & 15.2\% \\
\end{longtable}

The \(M_Z\) prediction of \(0.122\) matches the observed \(0.118\)
within 3\% (Figure 2). Figure 2 shows the model curve against the four
reference scales.

\begin{figure}
\centering
\pandocbounded{\includegraphics[keepaspectratio]{publication/figures/fig2_alpha_s_fix.png}}
\caption{α\_s from the φ⁴ associator layer model against QCD reference
values.}
\end{figure}

\subsubsection{3.8 Retrocausality, Baryogenesis, and Temporal
Consistency}\label{retrocausality-baryogenesis-and-temporal-consistency}

The probabilistic axiom of the directed-number algebra states that the
product \(0^0 \cdot 0^0\) yields a random variable drawn from a
distribution \(P(r)\). In IST we consider, as a working interpretation,
that this ``randomness'' is epistemic: \(r\) may be determined by global
boundary conditions that are not accessible to a local observer.

Operationally, the zero-point operator restricts allowable histories of
the substrate so as to enforce a global consistency condition. The term
``retrocausal'' is used only to describe the effective consequence that
constraints associated with future boundary data can influence present
conditional probabilities.

\textbf{Baryogenesis estimate:}

\[\eta \sim \frac{\alpha^4}{\varphi^2} \approx 1.1 \times 10^{-9}\]

which is within a factor of two of the observed value
\(\eta \approx 6 \times 10^{-10}\).

\paragraph{3.7.1 The Prime Number
Connection}\label{the-prime-number-connection}

The appearance of \(\alpha^{-1} \approx 137\) and \(\varphi\) in the
framework motivates a speculative organizing principle: stability may be
associated with arithmetic ``rigidity'' (e.g., primality for discrete
parameters) or with maximal incommensurability (e.g., strong
irrationality for continuous fixed points).

\begin{center}\rule{0.5\linewidth}{0.5pt}\end{center}

\subsection{4. Cosmological
Implications}\label{cosmological-implications}

\subsubsection{4.1 Overview}\label{overview-1}

This section sketches possible cosmological implications of a
non-orientable information substrate.

\subsubsection{4.2 The Flatness Problem}\label{the-flatness-problem}

The intrinsic geometry of a Klein bottle (\(K_2\)) is flat, so
Klein-bottle-based toy models provide a setting in which
\(\Omega_K \approx 0\) can arise without fine-tuning at the level of the
assumed topology. Within IST, the ``Big Bang'' is modeled as the
manifold's one-dimensional locus of self-intersection -- a structural
feature of the embedding used to motivate boundary conditions rather
than a detailed dynamical account of early-universe microphysics.

\subsubsection{4.3 Generalization to Higher
Dimensions}\label{generalization-to-higher-dimensions}

The framework is compatible with richer topologies, including
non-orientable quotients of Calabi-Yau manifolds. In such settings,
gauge symmetries (U(1), SU(2), SU(3)) may be associated with nontrivial
features of moduli space.

\subsubsection{4.4 Local Group Dynamics as a Consistency
Check}\label{local-group-dynamics-as-a-consistency-check}

Recent high-resolution constrained simulations {[}5{]} report that the
mass distribution around the Milky Way and Andromeda is strongly
flattened into a ``Local Sheet,'' with pronounced voids above and below.
In IST, this morphology is qualitatively consistent with an effectively
two-dimensional substrate whose coarse-grained projection yields a
three-dimensional volume. Sheet-like regions are associated with higher
inferred fold density and surrounding voids with lower fold density.

\subsubsection{4.5 Oscillatory Dark
Energy}\label{oscillatory-dark-energy}

The vacuum-pump cosmogony (Section 5.2) implies that the dark-energy
sector is not a static cosmological constant but a time-crystal
oscillation whose amplitude decays with redshift as
\(\varepsilon(z) = \varepsilon_0(1+z)^\beta\). A joint fit to 60
\(H(z)\) cosmic chronometers, 1701 Pantheon+ SNe Ia, and DESI DR1 BAO
compares the model against \(\Lambda\)CDM (Phase 16):

\begin{longtable}[]{@{}llll@{}}
\toprule\noalign{}
Model & \(\chi^2\) & \(\Delta\chi^2\) vs \(\Lambda\)CDM & \(H_0\) \\
\midrule\noalign{}
\endhead
\bottomrule\noalign{}
\endlastfoot
\(\Lambda\)CDM & 948 & -- & 73.6 \\
IST (\(\beta = 1/\varphi\)) & 926 & +22.1 & 71.4 \\
IST (free \(\beta\)) & 926 & +22.3 & 71.6 \\
\end{longtable}

The oscillatory model is preferred at approximately \(4\sigma\), and
\(H_0\) shifts from 73.6 to 71.4 km s⁻¹ Mpc⁻¹, pulling the Hubble
tension in the direction of the CMB-inferred value. The exponent was
tested across embedding dimensions (Phase 15c); \(d = 3\) is the clear
best fit, with the fitted \(\beta \approx 4.16\) within 2\% of the
associator volume prediction \(\beta = \varphi^3 = 4.236\) (Figure 3).

\begin{figure}
\centering
\pandocbounded{\includegraphics[keepaspectratio]{publication/figures/fig3_oscillatory_de.png}}
\caption{Oscillatory dark energy vs ΛCDM against H(z) data. The IST
model is preferred at \textasciitilde4σ (Δχ² = 22.1).}
\end{figure}

\subsubsection{\texorpdfstring{4.6 Void Lensing and the Suppression of
\(G\)}{4.6 Void Lensing and the Suppression of G}}\label{void-lensing-and-the-suppression-of-g}

The Phase 14 pinned \(G(\rho) \propto \rho^{1/\varphi}\) model was
applied to the void-lensing templates: low-density regions should show a
63\% suppression of the effective gravitational coupling, exactly the
golden-window prediction \(1 - (0.2)^{1/\varphi} = 0.63\) for a 10:1
density contrast. Forecasts with Euclid/COSMOS-Web depth distinguish
this from GR at \(10.7\sigma\) (Phase 17):

\begin{longtable}[]{@{}llll@{}}
\toprule\noalign{}
Model & Suppression & \(\chi^2\) vs GR & \(\sigma\) \\
\midrule\noalign{}
\endhead
\bottomrule\noalign{}
\endlastfoot
\(D = 2\) (grid) & 55.3\% & 88.2 & 9.4 \\
Phase 4 window & 61.9\% & 110.7 & 10.5 \\
\textbf{Phase 14 pinned} & \textbf{63.0\%} & \textbf{114.6} &
\textbf{10.7} \\
\end{longtable}

Real DES Y6 GOLD data produced a first stacked shear measurement from
3--4 voids with tangential shear \(\gamma_t \sim -0.025\) at \(0.27°\)
-- real signal but noise-limited at single-tile depth; a multi-tile
shear catalog is required for a decisive test.

\begin{center}\rule{0.5\linewidth}{0.5pt}\end{center}

\subsection{5. Simulation Results and Illustrative
Tests}\label{simulation-results-and-illustrative-tests}

This section reports the computational program (24 phases, 319 automated
tests; Python 3.14, numpy/scipy/numba, all reproducible) that
establishes the φ-attractor claim and tests the framework against
observation. The arc is: (i) falsify static φ (Phases 1--4); (ii)
demonstrate the anti-resonance attractor (Phases 6--9); (iii) establish
dynamical RG convergence and fold-density feedback (Phases 10--14); (iv)
close quantitative gaps (Phase 15); (v) test against real data (Phases
16--17); (vi) implement the plonk-scale substrate (Phases 23--24).

\subsubsection{5.1 Falsification of the Static-φ Hypothesis (Phases
1--4)}\label{falsification-of-the-static-ux3c6-hypothesis-phases-14}

\textbf{Phase 1 (Klein-bottle spectrum).} The substrate was modeled as a
discrete 4-regular twisted-torus graph cellulating the Klein bottle with
a flat \(\mathbb{Z}_2\) twist connection. The topological Laplacian was
verified (χ = 0, non-orientable, meridian holonomy −1) and the analytic
spectrum \(\lambda(p,\ell) = 4 - 2\cos(2\pi p/n) - 2\cos(\pi\ell/n)\)
validated to machine precision. \textbf{Two φ-tests failed:}
distinct-level gap ratios follow the number-theoretic \(4p^2 + \ell^2\)
ladder (median \(r^* \approx 0.77\)--0.92, no convergence to φ), and 2×2
block-spin RG preserves \(D_{\text{eff}} = 2\) with fixed point
\(D^* \approx 2\), not φ.

\textbf{Phase 2 (Hopf fibration).} A discrete Hopf fibration
\(S^1 \hookrightarrow S^3 \twoheadrightarrow S^2\) was constructed with
verified Chern number. The Kaluza--Klein relation \(\alpha = 4/R_f^2\)
with the topological minimum \(p = 3\) gives
\(\alpha_{\text{raw}} \approx 17.5\), far from
\(\alpha^{-1} \approx 137\); the required magnification
\(M \approx 49.0 \approx \varphi^8\) was identified but not derived at
this stage.

\textbf{Phase 3 (mass hierarchy).} The proton and electron formulas
remained at \(99.97\%\)/\(99.95\%\); the neutron was high by
\textasciitilde0.85 MeV and the naive associator \(\alpha_s\) model gave
0.38 vs observed 0.118 (both closed in Phase 15).

\textbf{Phase 4 (variable \(G\)).} Linearizing the Compression Operator
around the flat equilibrium gave the decay operator
\(M_\Psi = I - F^{-1}L/4\), with the slowest mode identified with the
gravitational time scale \(\tau_{\text{fold}} = 4/\gamma_{\min}\). A
fold-density scan showed \(D_{\text{eff}}\) descending from 3.43 to
1.17, \textbf{crossing φ exactly once at \(f \approx 4.20\)} -- the
first hint of the golden window, where the void suppression
\(1 - 1/f = 76.2\%\) matches IST phenomenology.

\textbf{Conclusion:} the local discrete topology is correct, but the
golden ratio is not present in the static graph. This is the required
negative control.

\subsubsection{5.2 The Anti-Resonance Attractor (Phases
6--9)}\label{the-anti-resonance-attractor-phases-69}

\textbf{Phase 6 (anti-resonance selection).} The golden rotation on the
spectral circle has a unique property: its gap rigidity
\(R = \min_{\text{gap}}/\max_{\text{gap}}\) stays at exactly
\(1/\varphi^2 \approx 0.382\) for all 300 simulated deposition
generations, while rational rotations \(p/q\) collapse exactly at
generation \(q + 1\). Non-noble irrationals survive at lower rigidity.
\textbf{The golden rotation is the unique maximal-persistence structure}
-- the same anti-resonance principle that produces Fibonacci spirals in
phyllotaxis. The Douady--Couder growth simulation converges to a
noble-family attractor at \(151.9° \pm 0.8°\). Because every finite
approximation is a Fibonacci rational converging to -- but never
reaching -- φ, the result is an attractor, not a fixed point.

\textbf{Phase 7 (vector substrate).} A non-raster ensemble of
oscillators on the spectral circle, coupled by spectral proximity: the
Fibonacci-golden ensemble is flat at \(D_{\text{eff}} = 1.10 \pm 0.03\)
across the 6--39° range, while random ensembles vary 0.5→2.2 and
rational (1/5) ensembles are chaotic and mode-locked.

\textbf{Phase 8 (vacuum-pump laser threshold).} The vacuum pump deposits
harmonic layers at golden-scaled positions \(f_k = f_0/\varphi^k\). A
sharp coherence transition occurs at \textbf{layer 11} -- the laser
threshold -- above which \(D_{\text{eff}}\) pins at 1.18 (the \(S^1\)
value), and the magnification at layer 8 matches \(\varphi^8 = 46.98\)
exactly. In 2D (Phase 8b), the spectral gap \(\lambda_{\min}\) grows
from \textasciitilde0 to 1.09 as golden accumulation activates
non-orientability.

\textbf{Phase 9 (golden phase selection).} A cellular automaton on the
Klein grid with golden-phase tracking shows the golden fraction rising
from 0.54→0.77 (+43\%) when live cells' phases rotate by the golden
angle per tick.

Figure 1 summarizes the interconnected mechanism: (A) anti-resonance
selection, (B) the vacuum-pump threshold, (C) dynamical RG convergence,
and (D) fold-density feedback.

\begin{figure}
\centering
\pandocbounded{\includegraphics[keepaspectratio]{publication/figures/fig1_phi_attractor_mechanism.png}}
\caption{The φ-attractor mechanism: (A) anti-resonance landscape with
maximal persistence at the golden rotation, (B) the vacuum-pump laser
threshold, (C) dynamical RG convergence near φ, (D) fold-density
feedback pinning at the golden window.}
\end{figure}

\subsubsection{5.3 Dynamical RG and the Golden Window (Phases
10--14)}\label{dynamical-rg-and-the-golden-window-phases-1014}

\textbf{Phase 12 (static RG fails).} Fibonacci-decimated blocking on the
golden-order circle produces nearly identical \(D_{\text{eff}}\) to
uniform blocking -- both far from φ. Static blocking of any kind cannot
converge to φ.

\textbf{Phase 13 (dynamical RG converges).} The blocking is not
pre-assigned; it emerges from the substrate's temporal evolution.
Golden-connected components (cells linked by edges with weight
\textgreater{} 0.5) become coarse vertices. Under the golden attractor,
\(D_{\text{eff}}\) \textbf{pins at \(1.655 \pm 0.001\)} from epoch 7
onward -- within 2.3\% of φ = 1.618.

\textbf{Phase 14 (fold-density feedback pins \(G_{\text{eff}}\)).} The
self-regulating ODE

\[\frac{df}{dt} = \gamma\,(D_{\text{eff}}(f) - \varphi)\,f\]

drives fold density to the golden window (\(f \approx 4.2\)) from any
initial condition, pinning the gravitational coupling at the exact
\(1/\varphi\) exponent. \(G_{\text{eff}} \propto \rho^{1/\varphi}\) is
therefore \textbf{derived}, not assumed.

\subsubsection{5.4 Closing the Quantitative Gaps (Phase
15)}\label{closing-the-quantitative-gaps-phase-15}

\begin{itemize}
\tightlist
\item
  \textbf{α\_s corrected} (Phase 15a): the associator layer count uses
  \(\varphi^4\) energy magnification per layer, yielding
  \(\alpha_s(M_Z) = 0.122\) (observed 0.118, 3\% error) and \(m_\tau\)
  within 1.3\% (Section 3.7, Figure 2).
\item
  \textbf{Neutron mass} (Phase 15b): running \(\varphi(\mu)\) reproduces
  \(m_n = 0.9395\) GeV (observed 0.9396, 99.99\%) (Section 3.6.7).
\item
  \textbf{Dimensional β} (Phase 15c): the redshift scaling of the DE
  amplitude is best fit by \(d = 3\), with \(\beta = \varphi^3 = 4.236\)
  within 2\% (Section 4.5).
\end{itemize}

\subsubsection{5.5 Observational Tests (Phases
16--17)}\label{observational-tests-phases-1617}

\begin{itemize}
\tightlist
\item
  \textbf{Oscillatory DE vs ΛCDM} (Phase 16): preferred at
  \textasciitilde4σ, \(\Delta\chi^2 = 22.1\), \(H_0\): 73.6→71.4
  (Section 4.5, Figure 3).
\item
  \textbf{Void lensing} (Phases 5, 17): the pinned \(G(\rho)\) model
  predicts 63\% suppression, distinguishable from GR at \(10.7\sigma\);
  real DES Y6 GOLD data produced a first stacked shear measurement
  \(\gamma_t \sim -0.025\) at \(0.27°\) (Section 4.6).
\end{itemize}

\subsubsection{5.6 The Plonk-Scale Substrate (Phases
23--24)}\label{the-plonk-scale-substrate-phases-2324}

The most recent phases implement a plonk-scale simulation with explicit
tracking of the 720° double-cover of the Klein bottle, addressing the
missing ingredient identified in Phases 19--22 (grid harmonics and gain
saturation corrupting the naive balloon models).

\textbf{Fibonacci lattice.} Golden-angle spiral positioning on the Klein
surface produces correlated phase-position ordering.

\textbf{4-tick orientation cycle.} Each oscillator advances one quarter
of the full Klein circumference per plonk tick; after 4 ticks (720°),
all oscillators return to their original chirality. The spin-1/2
double-cover is verified at 200/200.

\textbf{Parity-inverted coupling.} \(k_{\text{distance}}\) returns a
twist flag indicating whether the shortest geodesic crosses the Möbius
seam. 44.6\% of coupling entries are negative, encoding
orientation-reversing propagation. This prevents the uniform saturation
that plagued earlier balloon models and stabilizes amplitudes at
\textasciitilde0.91 (unsaturated, physically active). Figure 4 shows the
orientation cycle and the signed coupling structure.

\begin{figure}
\centering
\pandocbounded{\includegraphics[keepaspectratio]{publication/figures/fig4_parity_inversion.png}}
\caption{The plonk-scale substrate: (A) the Klein-bottle 720°
double-cover, (B) the 4-tick orientation cycle (chirality flips at each
twist crossing), (C) the parity-inverted signed coupling (44.6\%
negative), (D) stable-knot formation.}
\end{figure}

\textbf{Stable knots.} Knots form at a rate of \textasciitilde3\% per
4-tick cycle, robust across all parameter variations (Phase 24 sweep
over \(\omega_0\), gain, \(\sigma\), TOL, and \(N\)). Figure 5
demonstrates that the fraction is independent of the golden filter's
tolerance parameter -- the topological structure (Fibonacci lattice +
parity inversion) is the primary driver.

\begin{figure}
\centering
\pandocbounded{\includegraphics[keepaspectratio]{publication/figures/fig5_parameter_robustness.png}}
\caption{Stable-knot fraction (\textasciitilde3\%) is robust across the
parameter sweep.}
\end{figure}

\textbf{QM diagnostics} (Phase 23b): 100\% chirality flip at 180°
(spin-1/2), constructive/destructive superposition cycling, entanglement
via twist-geodesic pairs, and measurable phase-space uncertainty
(\(\Delta x \Delta p = 0.32\) vs the plonk bound 0.031).

\textbf{Scale bridging} (Phase 23c): plonk-scale knot formation maps to
the Compton and atomic scales via \(\varphi^8\) magnification (47×) and
golden-window \(G_{\text{eff}}\) pinning.

\textbf{Conclusion of the computational program.} The golden ratio acts
at the \textbf{structural level} -- Fibonacci lattice positions and
parity inversion through the Klein twist -- rather than as a tunable
filter parameter. φ emerges from how oscillators are positioned and how
they couple across the twist, validating the central claim: φ is a
dynamical attractor of the substrate's self-interaction, not an external
input.

\begin{center}\rule{0.5\linewidth}{0.5pt}\end{center}

\subsection{6. The Zero-Point Operator and Self-Referential Information
Automata}\label{the-zero-point-operator-and-self-referential-information-automata}

\subsubsection{6.1 Overview}\label{overview-2}

We define the Zero-Point Operator \(\Omega\), a specialization of the
Compression Operator \(\Psi\) that acts in regions of vanishing fold
density. The goal is twofold: (i) formalize \(\Omega\) as a geometric
limit of \(\Psi\), and (ii) describe how ``collapse/expansion'' dynamics
can be interpreted as a substrate-level mechanism for self-referential
information automata.

\subsubsection{\texorpdfstring{6.2 The Geometric Origin of \(\Psi\) and
\(\Omega\)}{6.2 The Geometric Origin of \textbackslash Psi and \textbackslash Omega}}\label{the-geometric-origin-of-psi-and-omega}

A recurring question: what performs the computation? If \(\Psi\) updates
the substrate state, what external agency applies the update? IST avoids
this regress by treating \(\Psi\) and \(\Omega\) as \textbf{geometric
necessities}: they express self-interaction in a non-orientable,
self-intersecting manifold. The ``computation'' is not performed by an
external entity; it is the manifold interacting with itself.

\subsubsection{6.3 Layered Substrates and the Directed-Number
Formalism}\label{layered-substrates-and-the-directed-number-formalism}

The operator formalism motivates a layered substrate (stacked sheets)
whose interference can generate an apparent 3D volume. This is developed
in the companion paper {[}4{]} on Directed Numbers and Zero-Point
Operators.

\subsubsection{6.4 The Plonk-Scale 720°
Double-Cover}\label{the-plonk-scale-720-double-cover}

Phase 23a implements the Zero-Point/Compression dynamics at the plonk
scale with explicit geometric content. Each oscillator carries a 4-state
orientation tracker encoding its position on the Klein bottle's
double-cover; one plonk tick advances the oscillator through one quarter
of the full Klein circumference. After 4 ticks -- a full 720° traversal
-- every oscillator returns to its original chirality, verifying the
spin-1/2 double-cover at 200/200 oscillators. This is the
substrate-level mechanism that underwrites fermionic spin: a parity
inversion (180°, one twist crossing) flips chirality, and only a second
full twist returns the state to itself.

The orientation cycle is implemented through parity-inverted coupling:
the shortest geodesic between two oscillators either stays on the same
sheet or crosses the Möbius seam, and the crossing is encoded as a sign
flip in the coupling matrix (44.6\% of entries negative). This
structural parity inversion is what prevents the uniform amplitude
saturation seen in earlier naive balloon models, stabilizing amplitudes
at \textasciitilde0.91 (unsaturated, physically active).

\subsubsection{6.5 Stable Knots and the Emergence of
Matter}\label{stable-knots-and-the-emergence-of-matter}

Stable knots -- localized soliton configurations in the substrate state
-- form at a rate of \textasciitilde3\% per 4-tick cycle, robust across
all parameter variations (Phase 24). The fraction is independent of the
golden filter's tolerance parameter, establishing that knot formation is
driven by topology (Fibonacci lattice positioning + parity inversion
through the twist), not by a tunable threshold. This is the
substrate-level counterpart of the topological solitons of Section 3.4.

\subsubsection{6.6 Quantum Diagnostics}\label{quantum-diagnostics}

Phase 23b verified at the plonk scale: 100\% chirality flip at 180°
(spin-1/2), constructive/destructive superposition cycling, entanglement
via twist-geodesic pairs, and measurable phase-space uncertainty
(\(\Delta x \Delta p = 0.32\) vs the plonk bound 0.031). These
diagnostics connect the substrate dynamics to quantum information
theory.

\subsubsection{6.7-6.9 {[}See companion paper for full
details{]}}\label{see-companion-paper-for-full-details}

\subsubsection{6.10 Implications for Consciousness and
Computation}\label{implications-for-consciousness-and-computation}

{[}Speculative section -- included as potential long-term direction, not
as scientific claim. See companion paper.{]}

\subsubsection{6.11 Toward a Generative
Ontology}\label{toward-a-generative-ontology}

{[}Philosophical implications -- work in progress.{]}

\begin{center}\rule{0.5\linewidth}{0.5pt}\end{center}

\subsection{7. Conclusion}\label{conclusion}

Information Substrate Theory provides a framework in which observed
physics emerges from a discrete, non-orientable two-dimensional
information substrate. A 24-phase computational program (319 automated
tests) establishes that the golden ratio φ is not a static invariant of
the substrate graph but a \textbf{dynamical attractor} of its harmonic
self-interaction -- emerging from anti-resonance selection, the
vacuum-pump laser threshold, and dynamical RG convergence, and pinned by
fold-density feedback.

\textbf{Key results:} - Proton mass:
\(M_P/m_p = (2/\varphi^2)\alpha^{-9}\) (99.966\% accuracy) - Strong
coupling: \(\alpha_s(M_Z) = 0.122\) from φ⁴ associator layers (observed
0.118, 3\% error) - Neutron mass: 0.9395 GeV with running φ(μ) (observed
0.9396, 99.99\%) - Variable gravity:
\(G_{\text{eff}} \propto \rho_{\text{fold}}^{0.618}\), derived via
fold-density feedback - Oscillatory dark energy: preferred over ΛCDM at
\textasciitilde4σ (\(\Delta\chi^2 = 22.1\)), \(H_0\): 73.6→71.4 - Void
lensing: 63\% suppression, distinguishable from GR at 10.7σ -
Plonk-scale: 720° double-cover verified (200/200), spin-1/2 flip at
100\%, stable knots \textasciitilde3\% robust

The framework is offered as a work in progress whose claims are intended
to be evaluated and revised against observation.

\begin{center}\rule{0.5\linewidth}{0.5pt}\end{center}

\subsection{Appendices}\label{appendices}

\subsubsection{Appendix A: Simulation
Protocols}\label{appendix-a-simulation-protocols}

{[}Detailed simulation procedures for reproducibility.{]}

\subsubsection{Appendix B: Reproducible Simulation
Code}\label{appendix-b-reproducible-simulation-code}

\paragraph{B.1 Code: CMB Parity
Analysis}\label{b.1-code-cmb-parity-analysis}

\begin{Shaded}
\begin{Highlighting}[]
\CommentTok{\# [Python code for CMB parity analysis]}
\end{Highlighting}
\end{Shaded}

\paragraph{B.2 Code: Void Gravity (Variable
G)}\label{b.2-code-void-gravity-variable-g}

\begin{Shaded}
\begin{Highlighting}[]
\CommentTok{\# [Python code for void gravity simulation]}
\end{Highlighting}
\end{Shaded}

\paragraph{B.3 Code: Topological Soliton
Generation}\label{b.3-code-topological-soliton-generation}

\begin{Shaded}
\begin{Highlighting}[]
\CommentTok{\# [Python code for soliton simulation]}
\end{Highlighting}
\end{Shaded}

\begin{center}\rule{0.5\linewidth}{0.5pt}\end{center}

\subsection{References}\label{references}

{[}1{]} NOWN Research Collective. Information Substrate Theory (IST): A
Topological Framework for Quantum Gravity and Consciousness. Internal
document v5.3 (2026).

{[}2{]} A. Solis. Golden Ratio as a Stable Infrared Fixed Point in
Nonlocal Field Theories. arXiv preprint (2025).

{[}3{]} Scognamiglio, D. et al.~COSMOS-Web: A Synthetic JWST Survey for
Void Lensing Anomalies. The Astrophysical Journal, in press (2026).

{[}4{]} NOWN Research Collective. Directed Numbers and Zero-Point
Operators: An Algebraic Formalism for Topological Information
Conservation. Internal document v0.7 (2026).

{[}5{]} Wempe, E., White, S.D.M., Helmi, A., Lavaux, G., \& Jasche, J.
The mass distribution in and around the Local Group. Nature Astronomy
(2026).

{[}6{]} van der Mark, M.B. \& 't Hooft, G.W. Light is Heavy.
arXiv:1508.06478v1 {[}physics.hist-ph{]} (2015).

{[}7{]} Williamson, J.G. \& van der Mark, M.B. Is the electron a photon
with toroidal topology? Annales de la Fondation Louis de Broglie, 22,
133 (1997).

{[}8{]} Oby, E.R., Degenhart, A.D., Grigsby, E.M., et al.~Dynamical
constraints on neural population activity. Nature Neuroscience, 28,
383-393 (2025).

{[}9{]} Tononi, G., Boly, M., Massimini, M., \& Koch, C. Integrated
information theory: from consciousness to its physical substrate. Nature
Reviews Neuroscience, 17, 450-461 (2016).

{[}10{]} Pletzer, B. et al.~The golden ratio in human brain dynamics.
Frontiers in Human Neuroscience, 9, 123 (2015).

{[}11{]} Khinchin, A. Ya. Continued Fractions. University of Chicago
Press (1964).

{[}12{]} Arnold, V. I. Proof of a theorem of A. N. Kolmogorov on the
preservation of conditionally periodic motions. Russian Mathematical
Surveys, 18(5), 9-36 (1963).

{[}13{]} CODATA Task Group on Fundamental Constants. 2022 CODATA
Recommended Values. Available at https://physics.nist.gov/cuu/Constants/
(2022).

{[}14{]} NOWN Research Collective. Decomposition of the Electron Mass
Factor \(12\pi^5\). Supplementary derivation, IST-workspace repository
(2026).

\end{document}
